\section{System Layout}

\subsection{Overview}

Phenopi is organized as a workflow rather than as a single application. The
main steps are image acquisition, image analysis, quality control, and growth
analysis. Each step reads input files from the experiment directory and writes
its own outputs, so intermediate results can be inspected before continuing. A
typical workflow starts with an experiment schedule. The Raspberry Pi uses this
schedule to capture images and store them in the image archive. The analysis
scripts then process these images, extract plant traits, and write the results
to tables. Diagnostic plots and overlays are used to check whether the images
and segmentations are suitable for downstream analysis.

\subsection{Hardware Components}

The acquisition setup consists of a Raspberry Pi, a Raspberry Pi Camera Module
v3, a fixed imaging frame, an experimental tray, and a stable light source. The
camera is mounted above the tray to acquire top-down RGB images. The hardware
setup should remain fixed during an experiment. Moving the camera, tray, or
light source can change the image geometry or illumination, which may affect
segmentation and trait extraction.

\subsection{Software Components}

Phenopi uses separate scripts for scheduling, image capture, image analysis,
plotting, and growth analysis. The scheduler runs on the Raspberry Pi and
starts capture jobs at the configured times. The capture script acquires
individual images and saves them using timestamped filenames. The analysis
scripts process the captured images after acquisition. They assign regions of
interest, segment plant canopies, extract traits, and save the results.
Plotting and growth analysis scripts use these trait tables to generate figures
and summary files.

\subsection{Directory Structure}

A typical Phenopi project or experiment directory contains folders for captured
images, runtime files, scripts, and results. The exact contents may differ
between installations, but the main folders should remain easy to identify.

\begin{forest}
  for tree={
    font=\ttfamily,
    grow'=0,
    child anchor=west,
    parent anchor=south,
    anchor=west,
    calign=first,
    edge path={
    \noexpand\path[draw, \forestoption{edge}]
    (!u.south west) + (7.5pt,0) |- (.child anchor)\forestoption{edge label};
    },
    before typesetting nodes={
    if n=1
    {insert before={[,phantom]}}
    {}
    },
    fit=band,
    before computing xy={l=15pt},
  }
  [phenopi/
    [captures/
      [capture\_20260508\_080000.jpg]
      [capture\_20260508\_080030.jpg]
      [\ldots]
    ]
    [deploy/
      [phenopi-scheduler.service]
    ]
    [results/]
    [runtime/
      [completed\_jobs.json]
      [schedule.json]
    ]
    [scripts/
      [analysis/]
      [capture/]
      [plotting/]
      [scheduling/]
    ]
  ]
\end{forest}

The most important configuration file for image acquisition is
\texttt{runtime/schedule.json}. This file defines when images should be
captured. The scheduler also writes its state to
\texttt{runtime/completed\_jobs.json}, which records which captures have
already been completed. Raw images are saved in \texttt{captures/} using
timestamped filenames. Analysis results are saved in \texttt{results/}, usually
as CSV files, diagnostic images, and plots. These files are the main outputs
used for quality control and biological interpretation. The \texttt{scripts/}
directory contains the command-line tools used for acquisition and analysis.
The \texttt{deploy/} directory contains the systemd service file used for
automated acquisition on the Raspberry Pi. 
